\section{Module Theory: Structure Theorem for Finitely Generated Modules over a PID}

\begin{theorem}[fg\_module\_structure]
Let $R$ be a principal ideal domain and $M$ a finitely generated
$R$-module. Then
\[
M \;\cong\; R^r \oplus R/(d_1) \oplus R/(d_2) \oplus \cdots \oplus R/(d_k)
\]
for some $r \geq 0$ and nonzero nonunits $d_1 \mid d_2 \mid \cdots \mid d_k$
in $R$, unique up to associates.
\end{theorem}

\begin{proof}
Present $M$ as $R^n / N$ for some submodule $N \leq R^n$, via a surjection
$R^n \twoheadrightarrow M$ from a finite generating set (finitely generated)
with kernel $N$. Since $R$ is a PID, $N$ is itself a free submodule of
$R^n$ of some rank $m \leq n$ (a submodule of a free module over a PID is
free). Choose bases of $R^n$ and $N$ so that the inclusion $N \hookrightarrow R^n$
is given by a diagonal matrix with entries $d_1 \mid d_2 \mid \cdots \mid d_m$
(the Smith normal form of the presentation matrix, obtained via elementary
row/column operations corresponding to change of basis on $R^n$ and $N$).
Then
\[
M = R^n / N \cong \bigoplus_{i=1}^m R/(d_i) \oplus R^{n-m},
\]
which is the stated decomposition with $r = n - m$, after discarding any
$d_i$ that are units (giving a zero summand). Uniqueness of the $d_i$ up to
associates follows from the invariance of the Smith normal form, which can
be characterized invariantly via the invariant factors
$\gcd$ of $i \times i$ minors of the presentation matrix.
\end{proof}
